\scp{Lexical items and features of Wallacea}{15-04}
While drafting this book,
we discussed how deeply to delve into the lexicon in this section,
and realised doing so in detail would require several more years
of work, and at least one additional volume.
So we chose to present a preliminary overview here,
and hold off on delving deeper into the lexicon for separate study.
While it would be both interesting
and insightful to track in more detail where lexical isoglosses
converge or overlap, we are also mindful that for decades the focus on the lexicon
has allowed the distractions that obscured insights facilitating
rigorous linguistic classification based on the historical phonology
in this region, and led scholars in wrong directions.
We are also mindful that highlighting the lexicon in this region
of intense mobility and contact could once again inappropriately elevate the role
of the lexicon in comparative studies in this region.
So we have made a decision to get this present study into the
hands of readers sooner,
and keep a more in depth study of the lexicon for a separate study.

Theoretically there are many kinds of things that can develop
in the lexicon, some of which are retentions,
some innovations, and some spread by diffusion or contact.
For example: 

\begin{itemize}
	\item	If a word exists in a proto-language,
				then we expect to find reflexes of that word in reasonably systematic
				ways (in form and meaning) in the branches
				and daughter languages that develop from that proto-language.
				For example, reflexes of PMP *matay ‘die’
				and *kutu ‘headlouse’ are fairly stable
				and widespread throughout Malayo-Polynesian languages and subgroups.
				Forms that reflect systematic sound correspondences
				and related meaning are seen as retentions (such as \it{mati,
				mata, mate}, all meaning ‘die’; and \it{kutu,
				koto, ʔutu, utu}, all meaning ‘headlouse’).
				Retentions are not useful for subgrouping.

	\item	If a new word in a subgroup replaces a term with the same meaning
				in the parent-language, then that innovation is a \it{lexical replacement}.
				For example, PMP *mata ‘eye’ is replaced by reflexes of **dama
				‘eye’ in all \tsc{Sula-Buru} languages,
				and reflexes are not found in other subgroups,
				except in the nearby Kayeli language due to contact with \tsc{Sula-Buru} languages.
				So **dama ‘eye’ is a lexical replacement assigned to Proto-Sula-Buru,
				and is taken as subgrouping evidence from the lexicon that supports
				the stronger evidence found in the historical phonology (see
				\trf{07-12}, and chapters 7 and 9).

	\item	A more subtle change is if a new form of an existing etymon
				replaces the previous form of the word in the parent-language,
				then that is another type of \it{lexical innovation}, or lexical development.
				So PMP *təbəl ‘thick’ underwent consonant metathesis
				and is fairly systematically reflected in several subgroups across
				central Maluku from a related form **batəlu ‘thick’.
				Reflexes of **batəlu do not seem to be found along the southern
				islands (where reflexes of PAn *kaS(ə)pal ‘thick’ are prevalent;
				see Grimes 1991a), but there is also no proto-language supported by the historical
				phonology to which to legitimately assign this etymon.
				\citet{Edwards2021} has many examples of this pattern for Proto-\tsc{Rote-Meto}
				which are assignable at various intermediate levels.
				In many cases such developments have subgrouping value,
				but in other cases they do not align with subgroups that are
				based on the historical phonology.

	\item	If the new form co-exists alongside a retention of the equivalent
				term from the parent-language,
				then that is an innovation by addition,
				but is not a replacement.
				For example, many languages across central Maluku have in one
				sense replaced PMP *ma-pənuq ‘full’,
				with reflexes of **ta(q)un ‘full’ for common uses of that meaning.
				However, reflexes of PMP *ma-pənuq ‘full’ are still found,
				albeit in more obscure uses (such as Buru \it{ben{\tl}beno} ‘sagging
				(from the weight of fruit)’).
				In many cases such developments may support subgroups that are
				based on stronger evidence,
				but in other cases they do not align with subgroups that are
				based on the historical phonology.

	\item	If the parent-language form is retained in some subgroups or
				nodes and a newer or different form is found sporadically
				and interspersed side-by-side in other subgroups or nodes,
				it is difficult to claim it as an innovation.
				Unless it aligns with other more convincing innovations,
				it has no real value for subgrouping.
				This uneven distribution of new forms interspersed with retentions
				is more likely to have been spread by \it{diffusion} or contact.
				In the face of criticisms regarding similar patterns of other
				features of PCMP and PCEMP, \citet[72]{Blust2009} conceded that such patterns
				should be seen as “a product of diffusion,
				since some languages that show no evidence of it are interspersed
				with those that do.” Much of the vocabulary assigned to PCMP
				and PCEMP falls into this latter category.
				Selected examples are given in the following sections.	
\end{itemize}

In \srf{15-04-01} we look at examples of words
and metaphors/phrases that are similar in the Wallacean
region that are found in both Austronesian and Papuan languages.
In these cases appealing to contact is unavoidable,
since these are not exclusive to Austronesian languages,
and several may have their ultimate source in Papuan forms.
In \srf{15-04-02} we look at examples of words that are not only
found in Wallacea, but are also found west of the Blust line,
such as in Sulawesi.
These are bigger than Blust's CMP or CEMP but many cannot be assigned to PMP.
In \srf{15-04-03} we look at examples of words
and features that are found in two or more Wallacean subgroups,
but not found in several others.
These are smaller than Blust's CMP and would not qualify for CEMP.
In \srf{15-04-04} we look at words
and features found in (parts of) our target region
and extending out to Oceania.
Some of these might be seen to support the CEMP hypothesis.
In \srf{15-04-05} we examine methodological implications for
our understanding of these different distributional ranges found
in the lexicon.